\documentclass[11pt]{article}
\usepackage[a4paper,margin=28mm]{geometry}
\usepackage{amsmath,amssymb,amsthm,mathtools,booktabs}
\usepackage[T1]{fontenc}\usepackage{lmodern,microtype}
\usepackage[colorlinks=true,linkcolor=blue,citecolor=blue,urlcolor=blue]{hyperref}
\newtheorem{theorem}{Theorem}[section]
\newtheorem{lemma}[theorem]{Lemma}
\newtheorem{proposition}[theorem]{Proposition}
\newtheorem{corollary}[theorem]{Corollary}
\theoremstyle{remark}\newtheorem{remark}[theorem]{Remark}
\newcommand{\Q}{\mathbb Q}\newcommand{\Z}{\mathbb Z}
\newcommand{\Fp}{\mathbb F_p}\newcommand{\Tr}{\operatorname{Tr}}
\title{Effective heights of elliptic sums with coprime indices}
\author{}\date{9 November 2026}
\begin{document}\maketitle
\begin{abstract}
Let an elliptic curve over a totally real number field have an integral
Weierstrass equation, and let $Q$ have infinite order.  For pairwise
coprime nonzero integer indices whose abscissas are positive in every
real embedding, the height of their sum is at least the maximum
individual height minus the finite-place height of $x(Q)$.  The loss is
independent of the number and size of the indices.  For the specific
curve over $\Q(\sqrt2)$ used in the explicit Diophantine constructions,
the bound is sharpened to $\min\{v_2(r),3\}\log2$ for $r$ summands;
for two coprime indices it is at most $\log2$.
Adding the fixed abscissa $x(P)$ when the number of summands is even
removes the loss entirely.
\end{abstract}

\section{The curve and height conventions}
Set $K=\Q(\rho)$, $\rho^2=2$, and
\[
d=3-\rho,\qquad D=d^2=11-6\rho,\qquad
E:\ y^2=x^3-Dx.
\]
The points
\[
R_0=(1+\rho,2\rho),\qquad
P=2R_0=\left(\frac{57-28\rho}{8},
                 \frac{-660+417\rho}{32}\right)
\]
belong to $E(K)$.  Denote its nontrivial Galois automorphism by
$\sigma$.  Write $x_n=x([n]P)$ for nonzero integers $n$.
We use the absolute logarithmic Weil height, with local valuations
normalized by $v(p)=1$ above a rational prime $p$.

The parameter $z=-x/y$ at the origin has
$v_2(z(P))=3/2$ at the unique place above $2$.
The formal logarithm converges there, has the same valuation as $z(P)$,
and is nonzero.  In particular, $P$ has infinite order.
At either real embedding, $x_n>d>1$: the point $[n]P$ is a nonzero
double, the duplication formula makes its abscissa positive, and the
curve equation then gives $x_n>d$.

\section{A height bound for coprime indices}
The basic argument applies more generally over any totally real number
field.  Write $h_{\mathrm{fin}}(z)$ for the sum of the finite-place
contributions to the absolute logarithmic height of $z$.

\begin{proposition}\label{prop:general}
Let $F$ be a totally real number field, let an elliptic curve over $F$
be given by an integral Weierstrass equation, and let $Q\in E(F)$ have
infinite order.  Suppose the abscissas $x([n_i]Q)$ are positive at every
real embedding, where the nonzero integers $n_1,\ldots,n_r$ are
pairwise coprime.  Then
\[
h\left(\sum_i x([n_i]Q)\right)
 \ge\max_i h(x([n_i]Q))-h_{\mathrm{fin}}(x(Q)).
\]
\end{proposition}

\begin{proof}
At a finite place where $Q$ is not formal, at most one $[n_i]Q$ is
formal, by Bezout's identity.  Hence at most one abscissa has a pole,
and no pole can cancel in the sum.
At a place above $p$ where $Q$ is formal, multiplication by an integer
prime to $p$ preserves the valuation of the formal parameter.  This
follows from $[n](z)=nz+O(z^2)$ with integral higher coefficients.
Multiplication by $p$ strictly increases its positive valuation.
At most one $n_i$ is divisible by $p$.  If such an index occurs, its
abscissa has strictly larger pole order than all the others, and
again no cancellation occurs.  Otherwise every summand has the same
local pole contribution as $x(Q)$, so the loss of that entire
contribution bounds any cancellation.  At each real embedding,
positivity makes the local height of the sum at least the local
height of each summand.  Summation proves the claim.
\end{proof}

We use the following elementary local properties; the full expansion
and its integrality at every finite place are proved in
\cite[Lemma~2]{EffectiveSums}.  For an integral
Weierstrass equation, the points with negative valuation of the
abscissa form the formal subgroup.  This remains true at a prime of bad
reduction.  On a sufficiently small formal neighborhood in characteristic
zero, its logarithm $\ell$ satisfies
\[
\ell([n]Q)=n\ell(Q),\qquad x(Q)-\ell(Q)^{-2}\in\mathcal O_v.
\]
For the equation above and the points generated by $P$ at $2$, these
identities follow directly from the formal expansions: $D$ is integral,
$v_2(z(P))=3/2$, and
\[
\ell(z)=\sum_{k\ge0}
 \frac{(-1)^k\binom{2k}{k}D^k}{4k+1}z^{4k+1}.
\]
The terms beyond the first have strictly larger valuation; expansion
of $x$ in $z$ shows the stated integral error.

\begin{theorem}\label{thm:coprime}
Let $n_1,\ldots,n_r$ be nonzero pairwise coprime integers and put
$S=\sum_{i=1}^r x_{n_i}$.  Then
\begin{equation}\label{eq:height}
h(S)\ge\max_i h(x_{n_i})-\min\{v_2(r),3\}\log2.
\end{equation}
If one index is even, the loss term may be omitted.  In particular,
for coprime nonzero $t,b$,
\[
h(x_t+x_b)\ge\max\{h(x_t),h(x_b)\}-\log2.
\]
\end{theorem}

\begin{proof}
Fix an odd finite place of $K$.  The coordinates of $P$ are integral
there, so $P$ is not in the formal subgroup.  If both $[n_i]P$ and
$[n_j]P$ belonged to that subgroup, Bezout's identity for coprime
$n_i,n_j$ would put $P$ in it.  Thus at most one summand has a pole.
If there is a pole, its valuation equals that of the sum.  Consequently
the local pole contribution of $S$ dominates that of each summand at
every odd finite place.

At the place above $2$, put $L=\ell(P)$, so $v_2(L)=3/2$.
The formal identities give
\[
x_{n_i}=n_i^{-2}L^{-2}+\epsilon_i,
\qquad \epsilon_i\in\mathcal O_v,
\qquad v_2(x_{n_i})=-3-2v_2(n_i).
\]
At most one index is even.  If one is even, its term has strictly
smaller valuation than all others and no cancellation occurs.
Otherwise all indices are odd, and
\[
L^2S=\sum_i n_i^{-2}+L^2\sum_i\epsilon_i,
\qquad
\sum_i n_i^{-2}\equiv r\pmod8.
\]
If $v_2(r)<3$, this gives $v_2(S)=-3+v_2(r)$.
If $8\mid r$, it gives $v_2(S)\ge0$.
The maximum local pole contribution of the summands is $3\log2$,
and its loss in $S$ is exactly
$\min\{v_2(r),3\}\log2$.  The local weight at this degree-two
completion is one in the absolute height.

At both real embeddings all summands exceed one.  Their sum therefore
has logarithmic local height at least that of each individual summand.
Summing the local inequalities proves \eqref{eq:height}.
\end{proof}

\begin{corollary}\label{cor:shift}
Under the hypotheses of Theorem~\ref{thm:coprime}, put
$\varepsilon=1$ when $r$ is even and $\varepsilon=0$ when $r$ is odd.
Then
\[
h\left(\sum_{i=1}^r x_{n_i}+\varepsilon x(P)\right)
 \ge\max_i h(x_{n_i}).
\]
In particular, for coprime nonzero $t,b$,
\[
h(x_t+x_b+x(P))\ge\max\{h(x_t),h(x_b)\}.
\]
\end{corollary}

\begin{proof}
If necessary, append the index $1$, which is coprime to every index.
The number of summands is then odd, so the loss in
Theorem~\ref{thm:coprime} is zero.
\end{proof}

\begin{remark}
For $r=4$ the loss is at most $2\log2$, in contrast to a loss depending
on the maximum index for arbitrary four-point sums.  Pairwise
coprimality is an additional hypothesis.  The source construction does
not impose it on all four root labels.
\end{remark}


\section{Application and scope}
If nonzero integers $a,t,b$ satisfy
\[
 t^2-a(a-2)tb+(a-2)^2b^2=1,
\]
then $\gcd(t,b)=1$, since every common divisor of $t$ and $b$ has its
square dividing the left side before the constant term.  Thus
Theorem~\ref{thm:coprime} applies to the two-point sum $x_t+x_b$.
This consequence is available only after the displayed integer
identity has been established, unless coprimality is separately
assumed or proved.

The theorem controls the height of an element of $K$.  It makes no
assertion about the height of that element's rational trace or norm.
Poles at distinct conjugate places can interact after those operations,
whereas the proof above compares each place separately.
\begin{thebibliography}{9}
\bibitem{EffectiveSums}
\emph{Effective heights of sums of elliptic abscissas},
project paper, 8 November 2026, Section~2, Lemma~2.
\href{https://github.com/JoseBrox/diophantine-z-q_improvements/blob/e9ee328a77f079b4869ea0cb61f60c686edba0ed/papers/effective_sums/effective_elliptic_sums_20261108_v10.tex}
{Source at the fixed project revision}.
\end{thebibliography}
\end{document}
